\documentclass[a4paper,12pt]{article}
\usepackage[utf8]{inputenc}
\usepackage[russian]{babel} 
\usepackage{graphicx} 
\usepackage{amsmath}
\usepackage{amssymb}
\usepackage{float}
\usepackage[colorlinks=true, linkcolor=blue]{hyperref}
\usepackage{enumitem}

\renewcommand{\floatpagefraction}{0.99}
\renewcommand{\topfraction}{0.99}      

\title{ Задача со стержнем.}
\author{Ларионов Станислав Александрович\and Доре Чойган Начын оглу \and Батанин Артём Александрович}
\date{\today}


\begin{document}

\maketitle

\newpage
\tableofcontents    
\newpage            


\section{Задача}
Объяснить физическую природу изменения слышимых частот вибраций алюминивого стержня



\section{Теоретическое предположение}

Предполагается, что длительное («протяжное») звучание стержня обусловлено возбуждением его собственных резонансных частот с малым коэффициентом затухания, причем высокочастотные гармоники затухают быстрее, чем главная, основная мода, что и способствует изменению тембра звука со временем.


\section{Физическая суть процесса}

В момент удара по стержню в нем возбуждается множество частот сразу. Энергия волн заперта внутри стержня и отражается от его торцов. Звук затухает медленно, потому что металл обладает малым внутренним трением. сама внутрення энерегия при возмущениях внутри самого стержня переходит в его нагрев.

\section{Оценка неопределённости измерений}
Неопределённость измерений вызвана следующими факторами:  дрейф температуры водно-ледяной смеси, тепловой шум в АЦП
\begin{itemize}[label={---}] 
    \item Дискретизация значений ЭДС в АПЦ;
    \item Внутреннее сопротивление АЦП и термопары;
    \item Дрейф температуры водно-ледяной смеси относительно 0°C;
    \item Tепловой шум в АЦП;
\end{itemize}
\subsection{Оценка систематической неопределённости}

\begin{itemize}[label={---}] 
    \item Дискретизация значений ЭДС в АЦП;
    Значение термо-ЭДС после оцифровки ограничены 5 знаками после запятой. Ниже приведен рассчёт неопределённости измерения ЭДС и температуры.
    \begin{equation} 
    \begin{aligned}
	\sigma_{E} = 10^{-5}/\sqrt{12}  \text{ В}\\
	\sigma_{E} \approx 2.89\cdot10^{-6} \text{ В}\\
	\sigma_{T} \approx 3.5\cdot10^{-5} \text{ °С}
    \end{aligned}
    \end{equation}
    \item Внутреннее сопротивление АЦП и термопары - $R \approx 10^7$ Ом $r \approx 10$ Ом 
    \begin{equation} 
    \begin{aligned}
	\sigma_{E} \approx R \cdot E/(R+r)  \text{ В}\\
	\sigma_{E} \approx 1.69\cdot10^{-8} \text{ В}\\
	\sigma_{T} \approx 2.0\cdot10^{-7} \text{ °С}
    \end{aligned}
    \end{equation}

    \item Дрейф температуры водно-ледяной смеси относительно 0°C: из-за низкой теплопроводности воды в стакане со смесью водя и льда мог создаться градиент температур от 0° до +24° (комнатная температура). На момент начала эксперимента градусник показал температуру смеси +8°C. Оценка неопределённости температуры: $\bar{T} = 4$ °C - ожидаемая температура  $\sigma_{T} = 4$ °C - дрейф температуры.
\end{itemize}

\section{Графики и результаты}

\section{Вклады авторов}

Доре Чойган - оформление LaTex

Ларионов Станислав - анализ физической природы

Батанин Артём - написание кода, анализ данных

Пальчиков - предоставил нам стержень для записи возмущений среды (звука)


\end{document}